\chapter{Contexte Général du Projet}
\label{chap:contexte_general}

\section{Présentation de l'Organisme d'Accueil}

\subsection{Description de l'Activité et des Produits/Services}
\textbf{Owl Media} est une agence digitale spécialisée dans le marketing, la communication et les solutions numériques innovantes. Elle propose une gamme complète de services permettant aux entreprises de développer leur présence sur le web et d'améliorer leur communication digitale.

Les principaux services de l'entreprise incluent :
\begin{itemize}
    \item \textbf{Création numérique et communication} : conception graphique, production vidéo, gestion des réseaux sociaux, campagnes publicitaires en ligne.
    \item \textbf{Stratégie digitale} : accompagnement sur les plans de communication, analyse de performance, SEO, marketing de contenu.
    \item \textbf{Solutions immersives Cameleon360°} : vidéos à 360°, visites virtuelles Matterport, intégration Google Street View.
\end{itemize}

\paragraph{Taille et effectif :}
\begin{itemize}
    \item L'entreprise compte plusieurs dizaines d'employés répartis dans différents départements spécialisés.
    \item Secteur d'activité : \textbf{marketing digital et production de contenu numérique}.
    \item La société opère principalement sur le marché marocain et régional (Afrique du Nord), mais certains services sont également proposés à l'international via le web.
\end{itemize}

\subsection{Historique et Positionnement}
Depuis sa création, \textbf{Owl Media} a su se démarquer par sa capacité à anticiper les tendances du marché et à innover dans ses méthodes de communication. Grâce à son savoir-faire éprouvé, elle est devenue au fil des années un partenaire stratégique pour de nombreuses entreprises souhaitant renforcer leur visibilité et développer leur image de marque sur le web. Aujourd'hui, l'agence se positionne comme le numéro 1 en Afrique du Nord dans son domaine, et son expérience accumulée constitue un véritable gage de confiance et de qualité.

\subsection{Organisation Interne et Organigramme}
L'organisation de Owl Media est structurée pour couvrir l'ensemble des besoins d'une agence digitale moderne. Chaque département a un rôle précis dans le fonctionnement global de l'entreprise.

\begin{figure}[htbp]
    \centering
    \includegraphics[width=0.9\textwidth]{organigramme_owlmedia.png}
    \caption{Organigramme structurel de l'entreprise Owl Media}
    \label{fig:organigramme_owlmedia}
\end{figure}

\paragraph{Organigramme commenté :}
Chaque service identifié sur la figure~\ref{fig:organigramme_owlmedia} remplit des fonctions clés :
\begin{enumerate}
    \item \textbf{Direction générale} : supervise l'ensemble des activités, prend les décisions stratégiques et représente l'entreprise auprès des clients et partenaires.
    \item \textbf{Service Ressources Humaines (RH)} : recrutement, formation, gestion du personnel et bien-être des employés.
    \item \textbf{Service Finance} : gestion de la comptabilité, budgets, paiements, suivi financier des projets.
    \item \textbf{Service Production / Création} : responsable de la conception graphique, production vidéo, animation et de la réalisation des projets clients.
    \item \textbf{Service Logistique} : planification et coordination des ressources nécessaires aux projets, gestion du matériel et des déplacements pour les tournages.
    \item \textbf{Service Système d'Information (SI)} : gestion des infrastructures informatiques, maintenance des serveurs, sécurité des données et support technique interne.
    \item \textbf{Service Qualité} : contrôle et validation des livrables, conformité avec les standards clients et internes, amélioration continue.
    \item \textbf{Service Maintenance} : entretien des équipements techniques, matériel informatique et studio, dépannage.
    \item \textbf{Service Ingénierie et Process} : innovation et développement de solutions techniques (ex. Cameleon360°), amélioration des processus de production et développement de solutions innovantes. Notre stage PFA s'inscrit au sein de ce pôle d'ingénierie et d'innovation IA.
\end{enumerate}

\subsection{Environnement de l'Entreprise : Clientèle et Fournisseurs}

\paragraph{Clientèle :}
Owl Media travaille principalement avec :
\begin{itemize}
    \item \textbf{Entreprises locales et régionales} : PME et grandes entreprises souhaitant améliorer leur présence digitale.
    \item \textbf{Institutions publiques et privées} : nécessitant des campagnes de communication et de marketing numériques.
    \item \textbf{Clients internationaux} : via les services web et les visites virtuelles.
\end{itemize}
\textbf{Origine des clients} : Maroc, Afrique du Nord, et quelques projets internationaux.\\
\textbf{Fréquence} : Les clients peuvent collaborer sur des projets ponctuels ou sur des contrats à long terme pour la gestion continue de leurs campagnes digitales.

\paragraph{Fournisseurs :}
Les fournisseurs de Owl Media sont essentiellement liés à :
\begin{itemize}
    \item \textbf{Matériel technique} : caméras, drones, ordinateurs, serveurs et autres équipements de production.
    \item \textbf{Logiciels et licences} : outils de création graphique, montage vidéo, réalité virtuelle et marketing digital.
    \item \textbf{Services externes} : prestataires pour la publicité en ligne, services cloud ou plateformes immersives.
\end{itemize}
\textbf{Origine des fournisseurs} : Maroc, Europe et parfois États-Unis (pour les logiciels spécialisés).\\
\textbf{Fréquence} : Les achats sont réalisés en fonction des besoins des projets et de l'entretien du matériel.

\section{Contexte et Motivation du Projet}

\subsection{Contexte Technologique : Évolution des LLM}
L'émergence des récents grands modèles de langage (\textit{Large Language Models} - LLM) tels que GPT-4 (OpenAI), LLaMA 3 (Meta) ou Mistral (Mistral AI) a révolutionné le domaine du traitement automatique du langage naturel (NLP). Ces modèles, entraînés sur des milliards de paramètres et des corpus textuels massifs, font preuve d'une capacité remarquable à synthétiser, reformuler et répondre à des questions complexes.

\subsection{Limites des LLM Classiques en Entreprise}
Malgré leurs performances exceptionnelles, l'utilisation directe des LLM commerciaux en milieu industriel se heurte à trois barrières fondamentales :
\begin{enumerate}
    \item \textbf{Le Phénomène d'Hallucination} : Les LLM ont tendance à générer des réponses plausibles mais factuellement fausses ou inventées lorsque l'information demandée ne fait pas partie de leur corpus d'entraînement.
    \item \textbf{Absence de Connaissances Privées et Récentes} : Les modèles génériques ne possèdent aucune connaissance des documents internes d'une entreprise (politiques RH, chartes informatiques, procédures internes).
    \item \textbf{Confidentialité et Souveraineté des Données} : L'envoi de documents internes sensibles vers des APIs cloud tierces est strictement proscrit par les réglementations et les politiques de sécurité d'Owl Media.
\end{enumerate}

\subsection{Émergence du RAG et Motivation d'Owl Media}
Pour surmonter ces limites, l'architecture \textbf{RAG} (\textit{Retrieval-Augmented Generation}) s'est imposée comme la solution idéale. Le RAG sépare le stockage des connaissances de la génération : il extrait d'abord les passages pertinents issus de la base documentaire d'entreprise via une recherche sémantique, puis injecte ces passages comme contexte exclusif au LLM pour la génération de la réponse.

Au sein d'Owl Media, la motivation d'un tel projet est manifeste : les collaborateurs perdent un temps précieux à rechercher manuellement des règles administratives ou des consignes techniques dans des fichiers PDF et Word dispersés. Développer la plateforme \textbf{Owl Media RAG MLOps} entièrement en local répond au besoin métier de fournir un assistant intelligent rapide, 100\% souverain et exempt d'hallucinations.

\section{Problématique et Enjeux}

\subsection{Formulation de la Problématique}
Le défi technique et opérationnel adressé dans ce mémoire s'énonce à travers la problématique suivante :

\begin{quote}
\textit{« Comment concevoir et déployer une plateforme RAG souveraine capable d'interroger les documents internes d'Owl Media en langage naturel tout en garantissant la fiabilité des réponses, l'absence d'hallucinations et une observabilité MLOps complète en production ? »}
\end{quote}

\subsection{Dimensions des Enjeux}
Cette problématique globale se décline selon quatre dimensions majeures :
\begin{itemize}
    \item \textbf{Enjeu Technique} : Dépasser la recherche vectorielle pure par l'association d'une recherche hybride (vectorielle dense + mots-clés BM25 sparse) et d'un réordonnancement par Cross-Encoder.
    \item \textbf{Enjeu Métier} : Offrir aux collaborateurs d'Owl Media un accès instantané en langage naturel aux informations RH et techniques, réduisant le temps de recherche documentaire de plusieurs dizaines de minutes à quelques secondes.
    \item \textbf{Enjeu de Confiance et Fiabilité} : Mettre en œuvre un mécanisme anti-hallucination rigoureux (évaluation du contexte et filtrage par seuil de score) garantissant que toute réponse s'appuie sur des sources authentiques.
    \item \textbf{Enjeu d'Observabilité MLOps} : Appliquer les standards d'ingénierie logicielle et MLOps aux systèmes IA (mesure des latences P50/P95, traçabilité par \texttt{query\_id}, traces sémantiques, et évaluation continue).
\end{itemize}

\section{Objectifs du Projet}

\subsection{Objectif Général}
L'objectif principal de ce projet est de développer et déployer la plateforme souveraine \textbf{Owl Media RAG MLOps}, une solution logicielle conteneurisée combinant un moteur de recherche hybride avancé, une génération locale supervisée et un tableau de bord d'observabilité MLOps en temps réel.

\subsection{Objectifs Spécifiques par Volet}
Les objectifs opérationnels se déclinent selon quatre volets techniques principaux :

\begin{enumerate}
    \item \textbf{Volet Ingestion et Vectorisation} :
    \begin{itemize}
        \item Extraire automatiquement le texte à partir de documents multiformats (PDF, DOCX, TXT).
        \item Découper les documents en fragments sémantiques (\textit{Semantic Chunking}).
        \item Générer les embeddings denses (via \texttt{nomic-embed-text}) et les indexer dans la base vectorielle locale Qdrant avec élimination des doublons.
    \end{itemize}

    \item \textbf{Volet Retrieval et Reranking} :
    \begin{itemize}
        \item Implémenter et comparer trois modes de recherche : Vectorielle sémantique, Mots-clés BM25 et Hybride.
        \item Intégrer un modèle Cross-Encoder (\textit{ms-marco-MiniLM-L-6-v2}) pour le reranking des résultats.
        \item Mettre en place un filtrage par seuil de pertinence pour prévenir les hallucinations.
    \end{itemize}

    \item \textbf{Volet Génération LLM Local} :
    \begin{itemize}
        \item Intégrer le modèle LLM local Llama 3.2 1B via Ollama pour la synthèse des réponses en langage naturel.
        \item Restituer de manière transparente les sources exactes (fichiers, extraits, scores de pertinence) associées à chaque réponse.
    \end{itemize}

    \item \textbf{Volet MLOps, Monitoring et Évaluation} :
    \begin{itemize}
        \item Enregistrer dynamiquement les métriques de chaque requête (latences par étape, statut, sources) dans une base SQLite dédiée (\texttt{monitoring.db}).
        \item Assurer la traçabilité sémantique fine via un système d'identifiant unique \texttt{query\_id} et l'intégration de traces distribuées (Langfuse).
        \item Automatiser l'évaluation quantitative des métriques RAG (\textit{Faithfulness}, \textit{Answer Relevance}, \textit{Answer Correctness}, \textit{Context Precision/Recall}, \textit{Hit Rate}) avec suivi MLflow et versionnement DVC.
        \item Proposer un tableau de bord web d'administration clair, sobre et réactif.
    \end{itemize}
\end{enumerate}

\section{Méthodologie de Travail Adoptée (Agile / Scrum)}

\subsection{Démarche Itérative et Rituels Scrum}
Afin de garantir une grande réactivité et un suivi rigoureux des livrables, nous avons adopté une méthodologie de développement Agile fondée sur les principes du cadre \textbf{Scrum}. Le projet s'est étalé sur une période de 2 mois, découpée en 5 Sprints itératifs de 1 à 2 semaines.

La gestion du projet a été rythmée par les activités suivantes :
\begin{itemize}
    \item \textbf{Sprint Planning} : Définition des objectifs et sélection des \textit{User Stories} du Product Backlog au début de chaque itération.
    \item \textbf{Points de Suivi Hebdomadaires} : Réunions régulières avec l'encadrant pour passer en revue l'avancement technique et ajuster la feuille de route.
    \item \textbf{Gestion des Versions et Intégration} : Utilisation rigoureuse de Git/GitHub pour le versionnement du code, de DVC pour la traçabilité des données de test, et de Docker Compose pour la conteneurisation des services.
\end{itemize}

\subsection{Planning et Découpage des Sprints}
La répartition des activités au cours des 5 Sprints du projet est synthétisée dans le tableau~\ref{tab:planning_sprints}.

\begin{table}[h!]
\centering
\caption{Planning récapitulatif des Sprints du projet PFA}
\label{tab:planning_sprints}
\begin{tabular}{|c|p{4cm}|p{8.5cm}|}
\hline
\textbf{Sprint} & \textbf{Intitulé} & \textbf{Livrables et Réalisations Principales} \\ \hline
\textbf{Sprint 1} & Ingestion \& Vectorisation & Configuration Docker, pipeline d'extraction multiformat, chunking sémantique, indexation dans Qdrant. \\ \hline
\textbf{Sprint 2} & Retrieval \& Recherche Hybride & Développement du moteur Vectoriel, Mots-clés BM25, fusion Hybride et intégration du Cross-Encoder Reranker. \\ \hline
\textbf{Sprint 3} & Génération LLM \& Anti-Hallucination & Connexion à Ollama (Llama 3.2 1B), ingénierie de prompt système, filtrage par seuil de pertinence. \\ \hline
\textbf{Sprint 4} & Couche MLOps \& Dashboard & Développement du service de monitoring SQLite, calcul des latences P50/P95, Live Logs, Trace Explorer, UI web épurée. \\ \hline
\textbf{Sprint 5} & Déploiement Docker, Évaluation \& Validation & Déploiement souverain multi-services (Docker Compose), automatisation des métriques RAG (Fidélité, Pertinence, Correctness), versionnement DVC, tracking MLflow et validation finale end-to-end. \\ \hline
\end{tabular}
\end{table}

\section{Conclusion}
Ce chapitre a permis de présenter le cadre général du projet, à savoir l'organisme d'accueil Owl Media, la motivation derrière le développement de la plateforme souveraine Owl Media RAG MLOps, ainsi que la problématique, les objectifs détaillés et la méthodologie de travail Agile/Scrum adoptée. 

Le chapitre suivant sera consacré à l'\textbf{État de l'art et étude théorique}, permettant de poser les bases théoriques indispensables concernant les modèles de langage, l'architecture RAG, les bases vectorielles et les principes du MLOps.
